\section{Formal Verification}
\label{sec:formal}

\subsection{Methodology: Refinement Mapping}

Directly model-checking the full system -- consensus, FSM, and network
sensors together -- is intractable at any useful scale. We instead specify a
layered hierarchy in TLA+~\cite{lamport2002specifying}: an abstract
consensus layer, a Tendermint-style consensus layer, an FSM layer, and a
server layer that composes them, and state that each concrete layer refines
the one above it. Liveness verification is reduced to a safety obligation on
a single linking invariant, \emph{QuorumOverlap} -- any two valid quorums
share at least one honest node -- following the LiDO framework's technique
for linearizable Byzantine distributed objects~\cite{qiu2024lido}, which
also requires the standard weak-fairness assumption on message delivery
after the partial-synchrony bound -- QuorumOverlap alone is a safety
invariant and does not by itself establish liveness without that fairness
assumption. The property this methodology targets is: under partial synchrony and an
$f$-bounded adversary, no two honest nodes ever decide conflicting blocks or
conflicting FSM states, and any safety violation implies the existence of
cryptographic double-signing evidence, provided (i) the health-aggregation
step is well-formed and quorum-certified, (ii) Bitcoin reorganizations are
no deeper than a configurable finality depth $K_{\text{DEEP\_FINALITY}}$,
and (iii) validators' peripheral-health views converge within the
partial-synchrony bound rather than being identical at every instant.
\textbf{We distinguish two levels of computer-checked evidence for this
property.} First, each layer's own local safety/liveness properties are
model-checked in isolation (Table~\ref{tab:formal-verification}); this
confirms each layer is internally free of the violations checked for,
within its own bounded configuration. Second, the refinement obligation
itself -- that the concrete, composed four-layer system actually simulates
the abstract layer above it, which is what would let the property transfer
from the abstract model to the concrete implementation -- is specified as
three configurations of increasing scale (C1: $N{=}4,f{=}1$; C2: deeper
consensus rounds at the same scale; C3: $N{=}7,f{=}2$) but \textbf{has not
reached convergence at any of the three scales, including the smallest
(C1), after a substantial run}. We do not report a completion estimate:
explicit-state model checking of a four-layer composition is subject to
combinatorial state-space growth, and an in-process run of this class of
check has continued well past $10^{8}$ generated states with no
discernible trend toward closure, on hardware modest enough that we cannot
responsibly project whether more compute would converge it in hours, days,
or not at all. We treat this as an open verification problem rather than a
pending task with a known cost, and return to it in
Section~\ref{sec:conclusion}. Table~\ref{tab:formal-verification}'s
per-layer results are therefore necessary but not yet sufficient evidence
for the property as stated; we report this distinction explicitly rather
than let the per-layer numbers stand in for the refinement result.

\subsection{Safety Invariants and Attack Resilience Lemmas}
\label{subsec:formal-invariants}

Five state invariants are specified and model-checked: \emph{Circuit Breaker
Isolation} (withdrawals are locked if and only if the FSM is in
\textsc{Sovereign} or \textsc{Recovering}); \emph{Hysteresis Integrity} (an
exit from \textsc{Recovering} to \textsc{Anchored} requires both the full
hysteresis wait and a valid re-anchoring proof); \emph{Suspicious-Exit
Hysteresis Integrity} (the symmetric guard on the
\textsc{Suspicious}$\to$\textsc{Anchored} edge); \emph{Strict FSM Transition
Safety} (only adjacent-state transitions are ever legal); and
\emph{Monotonicity Safety} (chain heights and time never move backward).
Beyond these state invariants, five attack-specific lemmas are enforced at
the proposal-validation layer: resistance to data withholding, resistance to
long-range/forged-anchor attacks via strict SPV verification, a
deterministic bound on Byzantine message flooding, resilience to eclipse
attacks via the six-bound P2P health check
(Section~\ref{subsec:sensors}), and accountability -- any safety violation
is designed to imply the existence of cryptographic double-signing evidence,
matching CometBFT's own equivocation evidence pool rather than requiring
custom slashing cryptography -- an accountability property inherited from
that evidence pool, not independently re-proven here.

Two threat-model boundaries -- restated here because they are direct
preconditions of the invariants above, not new assumptions -- are the
Bitcoin reorg-depth bound and the peripheral-view convergence assumption
introduced in Section~\ref{subsec:threat-model}.

\subsection{Verification Stress Test and Ablation Study}

\begin{table}[h]
\caption{Per-Layer, Non-Refinement Verification Results (TLC)}
\label{tab:formal-verification}
\centering
\footnotesize
\begin{tabular}{@{}lrrrr@{}}
\toprule
\textbf{Layer / property} & \textbf{States gen.} & \textbf{Distinct} & \textbf{Depth} & \textbf{Viol.} \\
\midrule
Tendermint (safety) & 3{,}573{,}347 & 509{,}883 & 25 & 0 \\
FSM (safety) & 1{,}022{,}325{,}986 & 3{,}386{,}761 & 22 & 0 \\
FSM (liveness) & 10{,}321{,}121 & 17{,}857 & 22 & 0 \\
\bottomrule
\end{tabular}
\end{table}

Table~\ref{tab:formal-verification} reports \emph{per-layer} model-checking
results -- each layer's own safety/liveness driver
run standalone, independent of the cross-layer refinement obligation
discussed above -- using each layer's own checked-in configuration
($N{=}4$, $f{=}1$ for Tendermint; a suspicious-state threshold of 2,
a sovereign-state threshold of 4, a hysteresis wait of 2 blocks, and a
down-hysteresis threshold of 2 for the FSM layer). All three runs
found zero violations across their explored, bounded state space. We note
this covers two of the four architectural layers described above (the
Tendermint-style and FSM layers); the abstract consensus layer and the
composing server layer do not have independent per-layer results of their
own outside the still-open refinement obligation below. As stated
above, this is evidence that each layer is individually well-behaved in
isolation; it does not by itself establish that the composed system
refines the abstract specification, which is what the still-open C1/C2/C3
sweep would show, for the reasons -- and with the same uncertainty about
cost and timeline -- discussed above; we report this honestly as an open
problem (Section~\ref{sec:conclusion}) rather than presenting the per-layer
numbers as a substitute.

The ablation study below spans two layers, not one: the circuit breaker,
hysteresis, and P2P health gate are FSM-layer mechanisms (the same scope
as the FSM rows in Table~\ref{tab:formal-verification}), while DA
consistency and the $f{+}1$ fast-forward are proposal-validation/consensus-layer
mechanisms. We state this explicitly because a mechanism's necessity is
shown only within its own layer's bounded configuration, not against the
full, composed system, until the refinement sweep above completes.

\begin{table}[h]
\caption{Ablation Study: Removing One Mechanism at a Time}
\label{tab:ablation}
\centering
\footnotesize
\begin{tabular}{@{}p{0.24\linewidth}p{0.24\linewidth}p{0.1\linewidth}p{0.24\linewidth}@{}}
\toprule
\textbf{Removed} & \textbf{Targeted threat} & \textbf{Depth} & \textbf{Violated property} \\
\midrule
Circuit breaker & Withdrawal leakage & 2 & \texttt{Sanity\_NeverAttempt} \texttt{WithdrawalLeakage} \\
Hysteresis & Premature re-anchoring & 6 & \texttt{Sanity\_NoIllegal} \texttt{HysteresisExit} \\
P2P health gate & False recovery under eclipse & 3 & \texttt{Sanity\_NoIllegal} \texttt{P2PRecovery} \\
DA consistency & Data withholding & 2 & \texttt{Sanity\_Never} \texttt{ProposeWithheldData} \\
$f{+}1$ fast-forward & Liveness deadlock & 4 & \texttt{EventualDecision} \texttt{UnderGSTLiveness} \\
\bottomrule
\end{tabular}
\end{table}

Table~\ref{tab:ablation} reports the ablation study: each row removes
exactly one defensive mechanism from the specification and re-runs
verification, expecting -- and finding -- a violation of a corresponding
sanity property. Four of five ablations were checked with Apalache's bounded
SMT search (each violation found within 25--120 seconds at a shallow depth
of 2--6), which explores depth-by-depth and is markedly faster than TLC's
breadth-first search for finding a shallow counterexample; the $f{+}1$
fast-forward ablation is liveness-only and was checked with TLC. Every
ablation reproduced a distinct, previously predicted failure mode with no
false negatives, which is the central evidence that each mechanism is
individually necessary to prevent the specific attack that motivated it
-- though single-factor ablation does not test interaction effects between
two or more mechanisms removed jointly, so this is evidence of individual
necessity, not evidence of joint sufficiency: removing the circuit
breaker permits exactly the withdrawal-leakage attack it exists to prevent;
removing hysteresis permits premature re-anchoring; removing the P2P health
gate permits an eclipsed node to force an illegal recovery; removing DA
consistency permits a data-withholding proposal to be accepted; and removing
the $f{+}1$ timeout fast-forward reintroduces a liveness deadlock under
Byzantine timeout manipulation.